% ch1.tex -- Module 8 ESSENTIALS edition, Chapter 1: Origins of the Elements (IQ1).
% book): same notation, same facts and numbers. Figures copied from src/mod8/figures/
% (same file names); the TikZ is copied (life-cycle chart made more compact) from the full
\chapter{Origins of the Elements: the Big Bang and the Stars}\label{ch1:chap}
\covers{M8.IQ1.01, M8.IQ1.02, M8.IQ1.03, M8.IQ1.04, M8.IQ1.05, M8.IQ1.06, M8.IQ1.07}

% nuclide symbol with right-aligned mass number A and atomic number Z: \chOneNuc{A}{Z}{X}
\newlength\chOneNucW
\newlength\chOneNucV
\newcommand\chOneNuc[3]{%
  \settowidth\chOneNucW{$\scriptstyle #1$}\settowidth\chOneNucV{$\scriptstyle #2$}%
  \ifdim\chOneNucV>\chOneNucW\chOneNucW=\chOneNucV\fi
  {}^{\makebox[\chOneNucW][r]{$\scriptstyle #1$}}_{\makebox[\chOneNucW][r]{$\scriptstyle #2$}}\mathrm{#3}}

\begin{summarybox}{In this chapter}
\begin{itemize}
  \item Energy into matter after the Big Bang; Hubble's evidence that space expands.
  \item Spectra and Wien's law; classes of stars; the HR diagram and the life of a star.
  \item Fusion in stars: the proton--proton chain, the CNO cycle and heavier elements.
\end{itemize}
\end{summarybox}

%% ==================================================================
\section{From radiation to matter: the Big Bang}\label{ch1:sec:bigbang}
\covers{M8.IQ1.01, M8.IQ1.03}

\begin{bigidea}
About 13.8 billion years ago the Universe was extremely hot and dense. As space expanded, it
cooled, and energy turned into matter ($E = mc^2$).
\end{bigidea}

\begin{keywords}
  \kw{Pair production}{a photon turning into a particle and its \textbf{antiparticle} (same
    mass, opposite charge).}
  \kw{Annihilation}{a particle and its antiparticle turning into photons.}
\end{keywords}

\begin{itemize}
  \item The \term{Big Bang} was not an explosion into empty space. Space itself stretches,
        everywhere at once, so there is no centre. Expanding space cools.
  \item $E = mc^2$ (on the formula sheet; $E$ in J, $m$ in kg): mass and energy are two forms
        of the same thing. Early high-energy \term{photons} (packets of light) with
        $E \geq 2mc^2$ made pairs (such as $\mathrm{e}^-$ and $\mathrm{e}^+$), near a nucleus or
        with a second photon. For a pair, use $2m$.
  \item As photons lost energy, proton pairs stopped first (they need more energy); electron
        pairs stopped a few seconds after the Big Bang. Matter and antimatter annihilated, but
        one particle in a billion was left: all matter today.
  \item Order: \term{inflation} (very fast expansion, \num{e-36} to \qty{e-32}{s}) $\to$
        \term{quarks} (tiny particles) join into protons and neutrons (\qty{e-6}{s}) $\to$
        helium nuclei form (3--20 minutes, \qty{e9}{K}).
  \item Then: atoms form and light is set free (\num{380000} years, \qty{3000}{K}) $\to$
        first stars (200 million years) $\to$ today (\qty{2.7}{K}).
  \item This \term{Big Bang nucleosynthesis} (making nuclei) left about \qty{75}{\percent}
        hydrogen and \qty{25}{\percent} helium by mass, and a trace of lithium. Heavier
        elements were made later, in stars.
\end{itemize}

%% ==================================================================
\section{Evidence for an expanding Universe}\label{ch1:sec:expansion}
\covers{M8.IQ1.02, M8.IQ1.01}

\begin{bigidea}
The further away a galaxy is, the faster it moves away (its light is redshifted). So space is
expanding, and long ago the Universe was small, hot and dense.
\end{bigidea}

\begin{itemize}
  \item Each element has its own set of \term{spectral lines} (dark lines in starlight). A source moving away stretches
        its waves, so all lines shift to longer wavelengths: \term{redshift} (towards us:
        \term{blueshift}). A bigger shift means a faster speed.
  \item \emph{Speeds}: Slipher (from 1912) found most galaxies were redshifted: moving away.
  \item \emph{Distances}: a \term{Cepheid variable} star brightens and dims in a regular cycle.
        Leavitt (1912): the cycle time gives its luminosity (true brightness); comparing this
        with how bright it looks gives its distance.
  \item Hubble used Cepheids to show Andromeda is another galaxy (1924). In 1929 he found
        $v \propto d$: twice as far, twice as fast. The gradient of $v$ against $d$ is the
        \term{Hubble constant} $H_0$, about \qty{70}{km.s^{-1}.Mpc^{-1}}.
\end{itemize}

% unbreakable: the card stays on one page (formula and symbols with their use and trap)
\begin{formulacard}[notonsheet, unbreakable]{Redshift, Hubble's law and the age of the Universe}
  \cardformula{z = \frac{\Delta\lambda}{\lambda_0} = \frac{\lambda - \lambda_0}{\lambda_0}
    \approx \frac{v}{c} \qquad v = H_0 d \qquad t \approx \frac{1}{H_0}}
  \begin{symbols}
    \sym{$z$; $\lambda_0$, $\lambda$}{redshift; laboratory and observed wavelength}{no unit; \si{m}}
    \sym{$v$; $d$}{recession velocity; distance ($\qty{1}{Mpc} = \qty{3.086e22}{m}$)}{\si{m.s^{-1}};
      \si{m}}
    \sym{$H_0$; $t$}{Hubble constant; age of the Universe}{\si{s^{-1}}; \si{s}}
  \end{symbols}
  \whentouse{$z \approx v/c$: light (not sound), $v$ below about $0.1c$. $v = H_0 d$: distant
    galaxies only (nearby ones have their own motion).}
  \derive{one wave every period $T$: $\lambda_0 = cT$, $\Delta\lambda = vT$, so
    $\frac{\Delta\lambda}{\lambda_0} = \frac{v}{c}$. At constant
    speed, $t = \frac{d}{v} = \frac{d}{H_0 d} = \frac{1}{H_0}$.}
  \trap{change $H_0$ from \si{km.s^{-1}.Mpc^{-1}} into \si{s^{-1}} for $t$: $H_0 = 70$ is
    \qty{2.27e-18}{s^{-1}}.}
\end{formulacard}

% Two more knowledge points of this section (in the full book: the subsections "What the
% expansion means" and "More evidence: the CMB and the light elements"), each with a run-in
% head as in "Classifying stars." below.
\textbf{What the expansion means.} The \term{cosmological principle}: on a large scale, the
Universe looks the same from everywhere. So every galaxy sees the others moving away: space
itself stretches, and Earth is not the centre. Run backwards, everything was once hot and dense.

\textbf{More evidence for the Big Bang.} The \term{cosmic microwave background} (CMB), found in
1965, is faint microwave radiation from every direction. It is the light set free at
\num{380000} years, stretched about 1100 times: today a \qty{2.7}{K} black body. Also, the
oldest stars have about \qty{25}{\percent} helium, as Big Bang nucleosynthesis predicts.

%% ==================================================================
\section{Spectra and the classification of stars}\label{ch1:sec:spectra}
\covers{M8.IQ1.04, M8.IQ1.05}

\begin{bigidea}
A hot, dense object gives a continuous spectrum set only by its temperature. A gas gives
bright or dark lines that identify the element. Dark lines sort stars into classes.
\end{bigidea}

\Cref{ch1:fig:spectra} shows a \term{continuous spectrum} (every wavelength), an
\term{emission spectrum} (bright lines only) and an \term{absorption spectrum} (dark lines).

\begin{figure}
  \centering
  \includegraphics{ch1_spectra.pdf}
  \caption{The three kinds of spectrum. Hydrogen's dark lines are at exactly the same
    wavelengths as its bright lines.}\label{ch1:fig:spectra}
\end{figure}

\begin{itemize}
  \item An electron in an atom can only have certain energies, its \term{energy levels}.
        \emph{Emission}: heat or a current lifts electrons; each drops back and gives out one
        photon, $E = hf$ = the gap.
  \item \emph{Absorption}: a photon with exactly the gap energy lifts an electron. It is given
        out again in all directions, so the line is dark. Same gaps, same wavelengths.
  \item A \term{black body} absorbs all light; its glow depends only on its temperature. Its
        particles jostle, so it gives out every wavelength. Hotter: more light at every
        wavelength, and a shorter peak (\textbf{Wien's displacement law}).
\end{itemize}

\begin{formulacard}{Wien's displacement law}
  \cardformula{\lambda_\mathrm{max} = \frac{b}{T}}
  \begin{symbols}
    \sym{$\lambda_\mathrm{max}$}{wavelength where the spectrum peaks}{\si{m}}
    \sym{$b$}{Wien's constant, \qty{2.898e-3}{m.K} (data sheet)}{\si{m.K}}
    \sym{$T$}{surface temperature}{\si{K}}
  \end{symbols}
  \whentouse{black bodies (stars, the CMB), not for emission or absorption lines.}
  \trap{$T$ in kelvin, not \si{\degreeCelsius}. $\lambda_\mathrm{max}$ comes out in metres.}
\end{formulacard}

\textbf{Classifying stars.} Starlight is an absorption spectrum: cooler outer gas absorbs some
wavelengths. The \term{spectral classes}, hottest to coolest, are O B A F G K M (``\textbf{O}h
\textbf{B}e \textbf{A} \textbf{F}ine \textbf{G}uy, \textbf{K}iss \textbf{M}e''), with
subclasses 0 to 9; the Sun is G2. (Ionised: an atom has lost electrons. ``Metal'': any element
heavier than helium.)

\begin{table}
  \caption{Spectral classes (temperatures are approximate). $\lambda_\mathrm{max}$ is from
    Wien's law.}\label{ch1:tab:classes}
  \centering\small
  \begin{tabular}{@{}C{9mm}L{27mm}L{20mm}L{17mm}L{58mm}@{}}
    \toprule
    \colhead{Class} & \colhead{Temperature (K)} & \colhead{Colour} &
      \colhead{$\lambda_\mathrm{max}$ (nm)} & \colhead{Main absorption lines}\\
    \midrule
    O & above \num{30000} & blue & below 97 & ionised helium; weak hydrogen\\
    B & \num{10000}--\num{30000} & blue-white & 97--290 & neutral helium; stronger hydrogen\\
    A & \num{7500}--\num{10000} & white & 290--386 & hydrogen strongest\\
    F & \num{6000}--\num{7500} & yellow-white & 386--483 & weaker hydrogen; ionised metals\\
    G & \num{5000}--\num{6000} & yellow & 483--580 & ionised calcium; many metals (Sun)\\
    K & \num{3500}--\num{5000} & orange & 580--828 & neutral metals; some molecules\\
    M & below \num{3500} & red & above 828 & molecules (titanium oxide); very weak hydrogen\\
    \bottomrule
  \end{tabular}
\end{table}

\begin{itemize}
  \item Payne (1925): the classes differ because of \emph{temperature}, not the elements.
        Hydrogen's visible lines need electrons in $n = 2$: strongest in class A. Cooler stars
        keep electrons in $n = 1$; in O stars hydrogen is ionised.
  \item A shift of all lines gives the star's motion towards or away from us.
  \item Line width gives gas density: broad lines mean a small, dense main sequence star;
        narrow lines a huge giant. This is the \term{luminosity class}: I supergiant, III
        giant, V main sequence (Sun: G2V).
  \item Fast rotation also broadens lines: one edge comes towards us, the other moves away.
\end{itemize}

%% ==================================================================
\section{The Hertzsprung--Russell diagram and the life of stars}\label{ch1:sec:hr}
\covers{M8.IQ1.06, M8.IQ1.07}

\begin{bigidea}
An HR diagram plots luminosity against surface temperature. A star's place shows its colour,
luminosity, size and stage of life. Its mass decides how it lives and dies.
\end{bigidea}

\term{Luminosity} ($L$) is the total power a star gives out (W); distance does not change it
($L_\odot = \qty{3.83e26}{W}$). \term{Apparent brightness}, how bright it looks, depends on
distance. \term{Absolute magnitude} ($M$) is how bright a star would look from 10 parsecs:
smaller is brighter, and 5 magnitudes is a factor of 100:
$\dfrac{L_A}{L_B} = 100^{(M_B - M_A)/5}$ (not on the formula sheet). On the HR diagram
(\cref{ch1:fig:hr}):
\begin{itemize}
  \item Hot stars are on the \textbf{left} (blue, class O); bright stars at the \textbf{top}.
  \item \term{Main sequence}: the diagonal band of stars fusing hydrogen into helium; about
        \qty{90}{\percent} of stars, because it is the longest stage.
  \item Top right: \term{red giants} and supergiants, cool but luminous, so huge. Bottom left:
        \term{white dwarfs}, hot but dim, so tiny. (Stars with the same temperature give out
        the same power from each square metre, so a dim hot star must be small.)
\end{itemize}

\begin{figure}
  \centering
  \includegraphics{ch1_hr.pdf}
  \caption{An HR diagram (simplified). The Sun: 1 = now, 2 = red giant, 3 = white
    dwarf.}\label{ch1:fig:hr}
\end{figure}

\textbf{The life of a star.} A gas cloud (\term{nebula}) collapses into a hot
\term{protostar}. At about 10~million~K hydrogen fusion starts: the main sequence, where
inward gravity balances outward gas pressure (\term{stellar equilibrium}). More mass ($M_\odot$ = the
Sun's mass) means faster fusion and a shorter life (Sun 10 billion years, $10\,M_\odot$ 30
million). Later it swells into a red giant or supergiant.

\begin{figure}
  \centering
  \begin{tikzpicture}[x=1mm, y=1mm,
      st/.style={concept=#1, font=\footnotesize, inner sep=1.4mm},
      fin/.style={concept=#1, font=\footnotesize, inner sep=1.2mm, fill=hscbrand!16}]
    \node[st=12mm] (neb) at (7.4,0) {nebula};
    \node[st=15mm] (pro) at (27.9,0) {protostar};
    \node[st=15mm] (ms) at (49.9,0) {main\\sequence};
    \node[st=17mm] (rg) at (76.9,9) {red giant};
    \node[st=17mm] (rsg) at (76.9,-9) {red\\supergiant};
    \node[st=15mm] (pn) at (99.9,9) {planetary\\nebula};
    \node[st=15mm] (sn) at (99.9,-9) {supernova};
    \node[fin=30mm] (wd) at (130,9) {white dwarf\\(core below $1.4\,M_\odot$)};
    \node[fin=30mm] (ns) at (130,-3) {neutron star\\(core $1.4$--$3\,M_\odot$)};
    \node[fin=30mm] (bh) at (130,-14) {black hole\\(core above $3\,M_\odot$)};
    \draw[link] (neb) -- (pro);
    \draw[link] (pro) -- (ms);
    \draw[link] (ms.east) -- (rg.west);
    \draw[link] (ms.east) -- (rsg.west);
    \node[font=\footnotesize\bfseries, text=hscbrand, anchor=south west] at (64,14.5)
      {Starting mass below $8\,M_\odot$};
    \node[font=\footnotesize\bfseries, text=hscbrand, anchor=north west] at (64,-14.5)
      {Starting mass above $8\,M_\odot$};
    \draw[link] (rg) -- (pn);
    \draw[link] (pn) -- (wd);
    \draw[link] (rsg) -- (sn);
    \draw[link] (sn.east) -- (ns.west);
    \draw[link] (sn.east) -- (bh.west);
  \end{tikzpicture}
  \caption{A star's life depends on its starting mass. A white dwarf is Earth-sized; a
    neutron star is about \qty{20}{km} across; light cannot escape a black
    hole.}\label{ch1:fig:lifecycle}
\end{figure}

%% ==================================================================
\section{Fusion in stars: \texorpdfstring{$E = mc^2$}{E = mc2} and nucleosynthesis}\label{ch1:sec:fusion}
\covers{M8.IQ1.03, M8.IQ1.07}

\begin{bigidea}
In \term{nuclear fusion}, light nuclei join. The products have a little less mass; this
\term{mass defect} $\Delta m$ (mass before minus mass after) becomes energy, $E = \Delta m\,c^2$.
\end{bigidea}

\begin{itemize}
  \item $4\,\chOneNuc{1}{1}{H} \to \chOneNuc{4}{2}{He}$ turns \qty{0.7}{\percent} of the mass
        into energy: \qty{6.4e14}{J} per kg, over 10 million times petrol. So stars shine for
        billions of years. Fusion releases energy only for nuclei up to iron.
  \item Fusion needs over 10 million K: nuclei are positive and repel, so only very fast ones
        get close enough to join.
\end{itemize}

\begin{formulacard}[combine={$E = mc^2$ and $\mathrm{u} = 931.5\ \mathrm{MeV}/c^2$}]{Energy released in a nuclear reaction}
  \cardformula{E = \Delta m\,c^2 \qquad E\,(\mathrm{MeV}) = \Delta m\,(\mathrm{u}) \times 931.5}
  \begin{symbols}
    \sym{$E$; $\Delta m$}{energy released; mass defect
      ($\qty{1}{u} = \qty{1.661e-27}{kg} = \qty{931.5}{MeV}/c^2$)}{\si{J} or \si{MeV}; \si{kg}
      or u}
  \end{symbols}
  \whentouse{known masses. Mass lost by a star of power $P$:
    $\frac{\Delta m}{\Delta t} = \frac{P}{c^2}$.}
  \derive{$E = \Delta m\,(\mathrm{u}) \times \frac{\qty{931.5}{MeV}}{c^2} \times c^2 =
    \Delta m\,(\mathrm{u}) \times \qty{931.5}{MeV}$.}
  \trap{one route only: u $\to$ MeV, or kg $\to$ J. Keep every decimal place.}
\end{formulacard}

\textbf{The proton--proton chain} is the main route in stars up to about $1.3\,M_\odot$ (core
below about 18 million K). It gives about \qty{26.7}{MeV} per helium-4, including the
positrons' annihilation ($\chOneNuc{2}{1}{H}$ =
deuterium, $\chOneNuc{0}{+1}{e}$ = positron, $\nu_\mathrm{e}$ = neutrino (neutral, almost no
mass); $A$ and $Z$ balance).

{\centering
\begin{tabular}{@{}l@{\qquad}l@{}}
  Step 1 (twice) & $\chOneNuc{1}{1}{H} + \chOneNuc{1}{1}{H} \to \chOneNuc{2}{1}{H}
    + \chOneNuc{0}{+1}{e} + \nu_\mathrm{e}$\\[2pt]
  Step 2 (twice) & $\chOneNuc{2}{1}{H} + \chOneNuc{1}{1}{H} \to \chOneNuc{3}{2}{He} + \gamma$\\[2pt]
  Step 3 & $\chOneNuc{3}{2}{He} + \chOneNuc{3}{2}{He} \to \chOneNuc{4}{2}{He} + \chOneNuc{1}{1}{H}
    + \chOneNuc{1}{1}{H}$\\[2pt]
  \textbf{Net} & $4\,\chOneNuc{1}{1}{H} \to \chOneNuc{4}{2}{He} + 2\,\chOneNuc{0}{+1}{e}
    + 2\nu_\mathrm{e} + 2\gamma$\\
\end{tabular}\par}

\textbf{The CNO cycle} is the main route in hotter stars (above about $1.3\,M_\odot$). Carbon-12
catches four
protons and is given back, so it is a \term{catalyst} (made by earlier stars). Net: the same
as the p--p chain, but with $3\gamma$.

{\centering\small
\begin{tabular}{@{}l@{\qquad}l@{}}
  1.\ $\chOneNuc{12}{6}{C} + \chOneNuc{1}{1}{H} \to \chOneNuc{13}{7}{N} + \gamma$ &
  4.\ $\chOneNuc{14}{7}{N} + \chOneNuc{1}{1}{H} \to \chOneNuc{15}{8}{O} + \gamma$\\[2pt]
  2.\ $\chOneNuc{13}{7}{N} \to \chOneNuc{13}{6}{C} + \chOneNuc{0}{+1}{e} + \nu_\mathrm{e}$ &
  5.\ $\chOneNuc{15}{8}{O} \to \chOneNuc{15}{7}{N} + \chOneNuc{0}{+1}{e} + \nu_\mathrm{e}$\\[2pt]
  3.\ $\chOneNuc{13}{6}{C} + \chOneNuc{1}{1}{H} \to \chOneNuc{14}{7}{N} + \gamma$ &
  6.\ $\chOneNuc{15}{7}{N} + \chOneNuc{1}{1}{H} \to \chOneNuc{12}{6}{C} + \chOneNuc{4}{2}{He}$\\
\end{tabular}\par}

\textbf{After the main sequence.} In a red giant (about 100 million K) the
\term{triple-alpha process} fuses $3\,\chOneNuc{4}{2}{He} \to \chOneNuc{12}{6}{C}$; then
$\chOneNuc{12}{6}{C} + \chOneNuc{4}{2}{He} \to \chOneNuc{16}{8}{O} + \gamma$. Massive stars fuse
carbon, neon, oxygen and silicon in shells, up to iron. Fusing iron takes in energy, so the
core collapses. Elements heavier than iron form by neutron capture (supernovae, neutron-star
collisions, old giant stars). Supernovae and planetary nebulae spread the elements into new stars.

%% ==================================================================
\section{Key experiments}\label{ch1:sec:experiments}
\covers{M8.IQ1.02, M8.IQ1.04, M8.IQ1.05, M8.IQ1.06}

\begin{experiment}{Hubble's law from galaxy data}
\textbf{Aim:} to show that $v \propto d$. \textbf{What was done:} Hubble (1929) found speeds
from redshifts and distances from Cepheids, then plotted $v$ against $d$. \textbf{Result:} a
straight line through the origin, so $v \propto d$. (Hubble's own gradient was about
\qty{500}{km.s^{-1}.Mpc^{-1}}: his distances were too small. Today
$H_0 \approx \qty{70}{km.s^{-1}.Mpc^{-1}}$, so $t \approx \frac{1}{H_0} = \qty{1.4e10}{years}$.)
\textbf{Conclusion and significance:} space is expanding: the first strong
evidence for the Big Bang (Hubble--Lema\^{\i}tre law).
\end{experiment}

\begin{experiment}{Observing continuous, emission and absorption spectra}
\textbf{Aim:} to compare the three kinds of spectrum. \textbf{What was done:} a spectroscope
(it splits light into its colours) viewed a filament lamp, gas discharge tubes, and the lamp through a sodium flame.
\textbf{Result:} a rainbow; bright lines for each gas (hydrogen: 656, 486, 434, \qty{410}{nm});
a dark sodium line at \qty{589}{nm}. \textbf{Conclusion and significance:} dark lines sit where
the bright lines are, so lines show what stars are made of.
\end{experiment}

\begin{experiment}{Plotting an HR diagram from star data}
\textbf{Aim:} to find the groups of stars. \textbf{What was done:} real stars plotted by
temperature (hot on the left) and absolute magnitude (bright at the top). \textbf{Result:}
the Sun, Sirius A and Vega on the main sequence; Betelgeuse and Rigel (supergiants) above it;
Sirius B (a white dwarf) below left. \textbf{Conclusion and significance:} stars group by size
and stage of life; the main sequence is the long, steady stage.
\end{experiment}

\begin{summarybox}{Remember}
\begin{itemize}
  \item Space expands and cools. Photons made pairs ($E \geq 2mc^2$); a little extra matter
        survived annihilation.
  \item Hubble: redshift + Cepheid distances give $v = H_0 d$. Also: the CMB and
        \qty{25}{\percent} helium.
  \item Continuous spectrum: set by $T$ ($\lambda_\mathrm{max} = \frac{b}{T}$). Lines: energy
        levels; dark and bright lines at the same wavelengths.
  \item O B A F G K M, hot to cool. HR: hot left, bright top. Below $8\,M_\odot$: white dwarf;
        above: supernova.
  \item $E = \Delta m\,c^2$. p--p chain or CNO cycle: 4 protons $\to$ helium-4,
        \qty{26.7}{MeV}. Heavier: up to iron in stars, then neutron capture.
\end{itemize}
\end{summarybox}
